% META: {"id":"1SPE-PROBCOND-EX-025","chapitre":"1SPE-PROBA-COND","type_objet":"exercice","capacites":["1SPE-PROBA-COND-C3"],"capacites_codes":["C3"],"methodes":["M3"],"parcours":1,"competences":["calculer","raisonner"],"duree_min":10,"mode_creation":"ex_nihilo","sources_inspiration":[],"parametres_sympy":{},"fichier_tex":"chapitres/1SPE-PROBA-COND/exercices/1SPE-PROBCOND-EX-025.tex","corrige_tex":"chapitres/1SPE-PROBA-COND/corriges/1SPE-PROBCOND-CO-025.tex","status":"approved","origin":"generated","status_history":["generated","audited","accepted","approved"],"release_acceptance":"RELEASE_OWNER_BATCH_ACCEPTANCE","acceptance_closure_digest":"sha256:9b3ccf9a81c5520fbb7e03b7b2d3e7bf2057a02834b6d908bf3be5f49f3b553f","acceptance_receipt_digest":"sha256:4fad86f0282e0fa4e0419cca1ad6379ae7af5a280c04a4a90880f90a1c044242"}
% BEGIN-VERIFY
% from sympy import Rational
% P_A = Rational(7, 10)
% P_B_A = Rational(4, 7)
% P_B_Abar = Rational(2, 3)
% P_B = P_A * P_B_A + (1-P_A) * P_B_Abar
% assert P_B == Rational(7,10)*Rational(4,7) + Rational(3,10)*Rational(2,3)
% assert P_B == Rational(2,5) + Rational(1,5)
% assert P_B == Rational(3, 5)
% END-VERIFY
\begin{exercice}{1SPE-PROBCOND-EX-025}{1}{10}
On donne $P(A) = \dfrac{7}{10}$, $P_A(B) = \dfrac{4}{7}$ et $P_{\overline{A}}(B) = \dfrac{2}{3}$. Calculer $P(B)$.
\end{exercice}
